\begin{proof}[Proof of Corollary~\ref{cor:quantitative-circuits}]
We use Williams's witness-circuit argument~\cite[Theorem~6.1]{Wil13}
with the local reduction of Jahanjou, Miles, and
Viola~\cite[Theorem~1 and Section~6]{JMV13}.

By \v{Z}\'ak's nondeterministic time hierarchy~\cite{Zak83}, in the
form stated in~\cite[Theorem~2.1 in the full version]{Wil13}, fix a unary language
$H\in\operatorname{NTIME}(2^n)\setminus\operatorname{NTIME}(2^{n/2})$
on multitape Turing machines. For all sufficiently large $n$ apply
the local reduction to its fixed machine on $1^n$ with locality parameter $r=2$. A time allowance
$T=2^n n^{O(1)}$ includes the polylogarithmic overhead of simulating
the machine in the reduction and gives $q=n+O(\log n)$ clause-index
bits. In polynomial time the reduction constructs a decoder giving
three variable labels and their signs on each index. Let $\varphi_n$
consist of the clauses decoded on all $2^q$ indices. It is satisfiable
exactly when $1^n\in H$.

Each of the decoder's $O(n)$ output bits has a constant-depth decision
tree, so it depends on at most a fixed number $w$ of index bits.
Evaluating the decoder on all strings of Hamming weight at most $w$
recovers each support. Indeed, two strings witnessing dependence on
one coordinate can be set to zero off the support while retaining
their differing outputs. The same evaluations give the truth table
on each support. This constructs a decoder of $O(n)$ fan-in-two gates
in polynomial time.

Pad each variable label to $dn$ bits for a fixed positive integer $d$
and declare all $2^{dn}$ variables. If $|i|=dn$ for such an $n$ and
$\varphi_n$ is satisfiable, let $W(i)$ be bit $i$ of its
lexicographically first satisfying assignment, ordering variables by
their labels. Set $W(i)=0$ in all other cases. The length of $i$
determines $n$. To decide $W$, write out the $2^q$ clauses, ask
an NP oracle whether the formula is satisfiable, and, if so,
recover the first assignment by asking whether its current prefix
followed by zero has a satisfying extension. At most $2^{dn}$ prefix
queries of length $2^{O(n)}$ suffice, so $W\in\mathrm E^{\mathrm{NP}}$.

Put $a(t)=\sqrt{\log_2 t}$. Testing whether the count in the
constant-density part of Theorem~\ref{thm:main} is positive
decides $3$-SAT with $V$ variables and at most $5V$ clauses in time
$V^b2^{A V/a(V)}$ for fixed $A,b>0$. We absorb the change of
logarithm base in $A$. The solver uses polynomial space, so a
multitape simulation can scan its workspace for each bit access at
a cost $V^{O(1)}$, absorbed in $b$.

Fix a positive rational $\eta<1/(12A)$ and put
$s_n=\lfloor\eta n\lfloor a(n)\rfloor\rfloor$ for $n\ge2$.
The threshold is computable using the binary length of $n$ and the
integer square root of $\lfloor\log_2 n\rfloor$. Suppose that
$C_W(dn)\le s_n$ for all sufficiently large $n$. To decide $H$,
reject nonunary inputs. On $1^n$, guess a circuit on $dn$ index bits
with at most $s_n$ gates. Feed the decoder's three variable labels to
three copies of this circuit and test whether the indexed clause is
violated. Accept if the test circuit is unsatisfiable. When
$1^n\in H$, a circuit computing $W$ at length $dn$ gives an
accepting guess. When $1^n\notin H$, every guess violates a clause.

The test circuit has $q$ inputs and at most $3s_n+O(n)$ gates. The
encoding above uses one variable per input and gate, at most four
clauses of width at most three per gate to enforce its truth table,
and one unit clause to pin the output. Padding with unused variables
gives $U_n=3s_n+Bn$ variables and at most $5U_n$ clauses for a fixed
integer $B$. Since $U_n=(3\eta+o(1))n a(n)$ and $a(U_n)\sim a(n)$,
the test takes $2^{(3A\eta+o(1))n}$ time. Guessing and building the
circuits take polynomial time, so $H$ is decidable nondeterministically
in time $2^{n/2}$ for all sufficiently large $n$. The finitely many
exceptional lengths can be hardwired, contradicting the choice of $H$.
Thus $C_W(dn)>s_n$ for infinitely many $n$. Since
$s_n\sim\eta n\sqrt{\log_2 n}$ and $d$ is fixed, taking $L=W$ and any
$0<c<\eta/(d\sqrt{\log 2})$ proves \eqref{eq:quantitative-circuit-lower-bound}.
\end{proof}
